\section{Extracción de características: de diseño manual a aprendizaje automático}

\subsection{Las limitaciones de la extracción manual de características}

Para clasificar imágenes correctamente, necesitamos extraer características que distingan las diferentes categorías. Los métodos tradicionales dependen de la \textbf{extracción manual de características} (Manual Feature Extraction), con el siguiente proceso:

\begin{figure}[H]
\centering
\includegraphics[width=0.9\textwidth]{slides-images/slide-15.jpg}
\caption{Proceso de extracción manual de características: conocimiento del dominio $\rightarrow$ definición de características $\rightarrow$ detección de características para la clasificación. El problema central es que este proceso es extremadamente sensible a los cambios.\protect\footnotemark}
\end{figure}
\footnotetext{Diapositiva 15.}

Por ejemplo, para reconocer un rostro, podrías definir las características como «nariz, ojos, boca». Pero esto plantea de inmediato un problema recursivo: ¿cómo detectar un «ojo»? Necesitarías definir características de nivel aún más bajo para describir la apariencia del ojo.

\begin{warningbox}{El dilema central de las características manuales}
La extracción manual de características enfrenta dos retos fundamentales:
\begin{enumerate}
    \item \textbf{La recursividad en la definición de características}: definir características de alto nivel requiere definir características de nivel más bajo, y la definición de estas últimas requiere características de un nivel aún más bajo, formando una recursión infinita.
    \item \textbf{La variabilidad}: la apariencia de un mismo tipo de objeto varía enormemente según el punto de vista, la iluminación, la escala, la deformación y la oclusión; las reglas hechas a mano no pueden enumerar todos los casos.
\end{enumerate}
\end{warningbox}

\begin{figure}[H]
\centering
\includegraphics[width=0.9\textwidth]{slides-images/slide-16.jpg}
\caption{Retos del reconocimiento de imágenes: cambios de punto de vista, cambios de escala, deformación, oclusión, condiciones de iluminación, fondo desordenado y variación intraclase.\protect\footnotemark}
\end{figure}
\footnotetext{Diapositiva 16.}

\subsection{Aprender jerarquías de características}

La idea central del aprendizaje profundo es: ¿es posible aprender directamente de los datos una jerarquía de características (hierarchy of features), en lugar de diseñarla a mano?

\begin{figure}[H]
\centering
\includegraphics[width=0.85\textwidth]{slides-images/slide-18.jpg}
\caption{Aprender la jerarquía de características a partir de los datos: características de bajo nivel (bordes, puntos oscuros) $\rightarrow$ características de nivel medio (ojos, orejas, nariz) $\rightarrow$ características de alto nivel (estructura facial).\protect\footnotemark}
\end{figure}
\footnotetext{Diapositiva 18. Referencia: Lee y colaboradores, ICML 2009.}

\begin{importantbox}{La idea clave de la jerarquía de características}
Las características visuales se organizan en capas:
\begin{itemize}
    \item \textbf{Características de bajo nivel} (Low-level features): bordes, manchas oscuras y otros elementos visuales básicos
    \item \textbf{Características de nivel medio} (Mid-level features): formadas por la combinación de características de bajo nivel, como ojos y orejas
    \item \textbf{Características de alto nivel} (High-level features): formadas por la combinación de características de nivel medio, como la estructura facial completa
\end{itemize}
El objetivo del aprendizaje profundo es que la red aprenda automáticamente este tipo de representación jerárquica de características.
\end{importantbox}

\subsection{Resumen de la sección}

La extracción manual de características no es escalable debido a los problemas de recursividad y variabilidad. La contribución central del aprendizaje profundo es permitir que la red aprenda automáticamente, a partir de los datos, una representación jerárquica de características, desde la detección de bordes de bajo nivel hasta la comprensión semántica de alto nivel, sin intervención humana en todo el proceso.

\newpage
